% Options for packages loaded elsewhere
\PassOptionsToPackage{unicode}{hyperref}
\PassOptionsToPackage{hyphens}{url}
\documentclass[
  11pt,
]{article}
\usepackage{xcolor}
\usepackage{amsmath,amssymb}
\setcounter{secnumdepth}{-\maxdimen} % remove section numbering
\usepackage{iftex}
\ifPDFTeX
  \usepackage[T1]{fontenc}
  \usepackage[utf8]{inputenc}
  \usepackage{textcomp} % provide euro and other symbols
\else % if luatex or xetex
  \usepackage{unicode-math} % this also loads fontspec
  \defaultfontfeatures{Scale=MatchLowercase}
  \defaultfontfeatures[\rmfamily]{Ligatures=TeX,Scale=1}
\fi
\usepackage{lmodern}
\ifPDFTeX\else
  % xetex/luatex font selection
\fi
% Use upquote if available, for straight quotes in verbatim environments
\IfFileExists{upquote.sty}{\usepackage{upquote}}{}
\IfFileExists{microtype.sty}{% use microtype if available
  \usepackage[]{microtype}
  \UseMicrotypeSet[protrusion]{basicmath} % disable protrusion for tt fonts
}{}
\makeatletter
\@ifundefined{KOMAClassName}{% if non-KOMA class
  \IfFileExists{parskip.sty}{%
    \usepackage{parskip}
  }{% else
    \setlength{\parindent}{0pt}
    \setlength{\parskip}{6pt plus 2pt minus 1pt}}
}{% if KOMA class
  \KOMAoptions{parskip=half}}
\makeatother
\usepackage{color}
\usepackage{fancyvrb}
\newcommand{\VerbBar}{|}
\newcommand{\VERB}{\Verb[commandchars=\\\{\}]}
\DefineVerbatimEnvironment{Highlighting}{Verbatim}{commandchars=\\\{\}}
% Add ',fontsize=\small' for more characters per line
\newenvironment{Shaded}{}{}
\newcommand{\AlertTok}[1]{\textcolor[rgb]{1.00,0.00,0.00}{\textbf{#1}}}
\newcommand{\AnnotationTok}[1]{\textcolor[rgb]{0.38,0.63,0.69}{\textbf{\textit{#1}}}}
\newcommand{\AttributeTok}[1]{\textcolor[rgb]{0.49,0.56,0.16}{#1}}
\newcommand{\BaseNTok}[1]{\textcolor[rgb]{0.25,0.63,0.44}{#1}}
\newcommand{\BuiltInTok}[1]{\textcolor[rgb]{0.00,0.50,0.00}{#1}}
\newcommand{\CharTok}[1]{\textcolor[rgb]{0.25,0.44,0.63}{#1}}
\newcommand{\CommentTok}[1]{\textcolor[rgb]{0.38,0.63,0.69}{\textit{#1}}}
\newcommand{\CommentVarTok}[1]{\textcolor[rgb]{0.38,0.63,0.69}{\textbf{\textit{#1}}}}
\newcommand{\ConstantTok}[1]{\textcolor[rgb]{0.53,0.00,0.00}{#1}}
\newcommand{\ControlFlowTok}[1]{\textcolor[rgb]{0.00,0.44,0.13}{\textbf{#1}}}
\newcommand{\DataTypeTok}[1]{\textcolor[rgb]{0.56,0.13,0.00}{#1}}
\newcommand{\DecValTok}[1]{\textcolor[rgb]{0.25,0.63,0.44}{#1}}
\newcommand{\DocumentationTok}[1]{\textcolor[rgb]{0.73,0.13,0.13}{\textit{#1}}}
\newcommand{\ErrorTok}[1]{\textcolor[rgb]{1.00,0.00,0.00}{\textbf{#1}}}
\newcommand{\ExtensionTok}[1]{#1}
\newcommand{\FloatTok}[1]{\textcolor[rgb]{0.25,0.63,0.44}{#1}}
\newcommand{\FunctionTok}[1]{\textcolor[rgb]{0.02,0.16,0.49}{#1}}
\newcommand{\ImportTok}[1]{\textcolor[rgb]{0.00,0.50,0.00}{\textbf{#1}}}
\newcommand{\InformationTok}[1]{\textcolor[rgb]{0.38,0.63,0.69}{\textbf{\textit{#1}}}}
\newcommand{\KeywordTok}[1]{\textcolor[rgb]{0.00,0.44,0.13}{\textbf{#1}}}
\newcommand{\NormalTok}[1]{#1}
\newcommand{\OperatorTok}[1]{\textcolor[rgb]{0.40,0.40,0.40}{#1}}
\newcommand{\OtherTok}[1]{\textcolor[rgb]{0.00,0.44,0.13}{#1}}
\newcommand{\PreprocessorTok}[1]{\textcolor[rgb]{0.74,0.48,0.00}{#1}}
\newcommand{\RegionMarkerTok}[1]{#1}
\newcommand{\SpecialCharTok}[1]{\textcolor[rgb]{0.25,0.44,0.63}{#1}}
\newcommand{\SpecialStringTok}[1]{\textcolor[rgb]{0.73,0.40,0.53}{#1}}
\newcommand{\StringTok}[1]{\textcolor[rgb]{0.25,0.44,0.63}{#1}}
\newcommand{\VariableTok}[1]{\textcolor[rgb]{0.10,0.09,0.49}{#1}}
\newcommand{\VerbatimStringTok}[1]{\textcolor[rgb]{0.25,0.44,0.63}{#1}}
\newcommand{\WarningTok}[1]{\textcolor[rgb]{0.38,0.63,0.69}{\textbf{\textit{#1}}}}
\usepackage{longtable,booktabs,array}
\newcounter{none} % for unnumbered tables
\usepackage{calc} % for calculating minipage widths
% Correct order of tables after \paragraph or \subparagraph
\usepackage{etoolbox}
\makeatletter
\patchcmd\longtable{\par}{\if@noskipsec\mbox{}\fi\par}{}{}
\makeatother
% Allow footnotes in longtable head/foot
\IfFileExists{footnotehyper.sty}{\usepackage{footnotehyper}}{\usepackage{footnote}}
\makesavenoteenv{longtable}
\setlength{\emergencystretch}{3em} % prevent overfull lines
\providecommand{\tightlist}{%
  \setlength{\itemsep}{0pt}\setlength{\parskip}{0pt}}
% Shared preamble for the LAMMPS-CO manuscript, supplementary, and cover letter.
% Compiled with xelatex. Latin Modern is used rather than a system font so the
% build is reproducible on any TeX installation; the characters it lacks
% (Greek, super/subscripts, arrows, a few symbols) are mapped explicitly below.
\usepackage{fontspec}
\setmainfont{Latin Modern Roman}
\setsansfont{Latin Modern Sans}
\setmonofont{Latin Modern Mono}[Scale=0.85]

\usepackage[a4paper,margin=1in]{geometry}
\usepackage{longtable,booktabs,array}
\usepackage{graphicx}
\usepackage{amsmath,amssymb}
\usepackage[table]{xcolor}
\usepackage[hidelinks,breaklinks=true]{hyperref}
\usepackage{microtype}
\usepackage{newunicodechar}

% Characters that appear in the text but not in Latin Modern's text font.
\newunicodechar{λ}{\ensuremath{\lambda}}
\newunicodechar{Δ}{\ensuremath{\Delta}}
\newunicodechar{ε}{\ensuremath{\varepsilon}}
\newunicodechar{σ}{\ensuremath{\sigma}}
\newunicodechar{ρ}{\ensuremath{\rho}}
\newunicodechar{τ}{\ensuremath{\tau}}
\newunicodechar{φ}{\ensuremath{\varphi}}
\newunicodechar{δ}{\ensuremath{\delta}}
\newunicodechar{α}{\ensuremath{\alpha}}
\newunicodechar{θ}{\ensuremath{\theta}}
\newunicodechar{→}{\ensuremath{\rightarrow}}
\newunicodechar{≈}{\ensuremath{\approx}}
\newunicodechar{√}{\ensuremath{\sqrt{}}}
\newunicodechar{∫}{\ensuremath{\int}}
\newunicodechar{⟨}{\ensuremath{\langle}}
\newunicodechar{⟩}{\ensuremath{\rangle}}
\newunicodechar{■}{\ensuremath{\blacksquare}}
\newunicodechar{⁺}{\ensuremath{^{+}}}
\newunicodechar{⁻}{\ensuremath{^{-}}}
\newunicodechar{⁰}{\ensuremath{^{0}}}
\newunicodechar{⁴}{\ensuremath{^{4}}}
\newunicodechar{⁵}{\ensuremath{^{5}}}
\newunicodechar{⁶}{\ensuremath{^{6}}}
\newunicodechar{⁷}{\ensuremath{^{7}}}
\newunicodechar{₀}{\ensuremath{_{0}}}
\newunicodechar{₁}{\ensuremath{_{1}}}
\newunicodechar{₂}{\ensuremath{_{2}}}
\newunicodechar{₃}{\ensuremath{_{3}}}
\newunicodechar{₄}{\ensuremath{_{4}}}
\newunicodechar{₆}{\ensuremath{_{6}}}

% pandoc emits this when a document mixes tables and lists
\providecommand{\tightlist}{%
  \setlength{\itemsep}{0pt}\setlength{\parskip}{0pt}}

% Long tables of file names need to break across pages
\setlength{\LTleft}{0pt}
\setlength{\LTright}{0pt}
\renewcommand{\arraystretch}{1.15}

\setlength{\parskip}{0.6em}
\setlength{\parindent}{0pt}
\sloppy
\usepackage{bookmark}
\IfFileExists{xurl.sty}{\usepackage{xurl}}{} % add URL line breaks if available
\urlstyle{same}
% fallback for those not using the hyperref driver hyperxmp:
\makeatletter
\@ifundefined{xmpquote}{\newcommand{\xmpquote}[1]{#1}}{}
\makeatother
\hypersetup{
  hidelinks,
  pdfcreator={LaTeX via pandoc}}

\author{}
\date{}

\begin{document}

\section{Supplementary Materials}\label{supplementary-materials}

\subsection{AI-Assisted Systematic Optimization of the LAMMPS Molecular
Dynamics Simulator via Large Language Model Coding
Agents}\label{ai-assisted-systematic-optimization-of-the-lammps-molecular-dynamics-simulator-via-large-language-model-coding-agents}

\textbf{Chun-Yeol You}

\begin{center}\rule{0.5\linewidth}{0.5pt}\end{center}

\subsection{Contents}\label{contents}

\begin{itemize}
\tightlist
\item
  Supplementary Table S1: Complete list of 35 new OMP pair variants
\item
  Supplementary Table S2: A-3 restrict/fxtmp modification log (102
  files, representative entries)
\item
  Supplementary Table S3: OPLS-AA force field parameters used in FEP
  Case Study
\item
  Supplementary Table S4: Li⁺ Joung--Cheatham parameters and
  Lorentz--Berthelot mixed coefficients
\item
  Supplementary Table S5: OPLS-AA FEP dU/dλ by lambda window (all
  solvents)
\item
  Supplementary Table S6: A-3 restrict/fxtmp optimization scope by
  category
\item
  Supplementary Table S7: A-3 before/after performance benchmarks (three
  representative styles)
\item
  Supplementary Methods S1: Claude Code prompt templates (OMP port
  generation)
\item
  Supplementary Methods S2: OPLS-AA system construction and FEP protocol
\item
  Supplementary Figure S1: Full OMP scaling benchmarks for new styles
\item
  Supplementary Figure S2: NVE energy conservation comparison
\item
  Supplementary Figure S3: A-3 code transformation (before/after pair
  hotloop)
\end{itemize}

\begin{center}\rule{0.5\linewidth}{0.5pt}\end{center}

\subsection{Supplementary Table S1: Complete List of 35 New OMP Pair
Variants}\label{supplementary-table-s1-complete-list-of-35-new-omp-pair-variants}

All styles below were generated by Claude Code, compiled with zero new
MSVC warnings, and verified bit-identical to their serial counterparts
by a 50-step NVE thermo comparison
(\texttt{thermo\_modify\ format\ float\ \%20.15g}). Files are located in
\texttt{src/OPENMP/} with the naming convention
\texttt{pair\_{[}style\_name{]}\_omp.\{h,cpp\}}.

% Table S1 rendered in landscape so wide \nolinkurl{} style names fit without overflow
\begin{landscape}
\begingroup
\small
\setlength{\tabcolsep}{5pt}
\renewcommand{\arraystretch}{1.18}
% landscape text width = 11in - 2in margins = 9in = 22.86cm
% column allocation: r(0.4) + 2.4 + 4.8 + 5.4 + 2.8 + 3.2 = 19.0 + sep(1.8) = 20.8cm < 22.86cm
\begin{longtable}{r l p{4.8cm} p{5.4cm} l p{3.2cm}}
\toprule
\# & Package & Standard style & OMP variant & Loop type & New register \\
\midrule
\endfirsthead
\multicolumn{6}{l}{\small\itshape (Table S1 continued\ldots)}\\[2pt]
\toprule
\# & Package & Standard style & OMP variant & Loop type & New register \\
\midrule
\endhead
\midrule
\multicolumn{6}{r}{\small\itshape (continued on next page)}\\
\endfoot
\bottomrule
\endlastfoot
 1 & EXTRA-PAIR & \nolinkurl{lj/cut/soft}                  & \nolinkurl{lj/cut/soft/omp}                  & half-Newton      & ---                  \\
 2 & EXTRA-PAIR & \nolinkurl{lj/cut/coul/long/soft}         & \nolinkurl{lj/cut/coul/long/soft/omp}         & half-Newton      & ---                  \\
 3 & EXTRA-PAIR & \nolinkurl{lj/class2/soft}                & \nolinkurl{lj/class2/soft/omp}                & half-Newton      & ---                  \\
 4 & EXTRA-PAIR & \nolinkurl{lj/cut/coul/cut/soft}          & \nolinkurl{lj/cut/coul/cut/soft/omp}          & half-Newton      & ---                  \\
 5 & EXTRA-PAIR & \nolinkurl{lj/cut/coul/dsf/soft}          & \nolinkurl{lj/cut/coul/dsf/soft/omp}          & half-Newton      & ---                  \\
 6 & EXTRA-PAIR & \nolinkurl{morse/soft}                    & \nolinkurl{morse/soft/omp}                    & half-Newton      & ---                  \\
 7 & EXTRA-PAIR & \nolinkurl{coul/cut/soft/gapsys}          & \nolinkurl{coul/cut/soft/gapsys/omp}          & half-Newton      & ---                  \\
 8 & EXTRA-PAIR & \nolinkurl{nm/cut/split}                  & \nolinkurl{nm/cut/split/omp}                  & half-Newton      & ---                  \\
 9 & EXTRA-PAIR & \nolinkurl{born/coul/dsf}                 & \nolinkurl{born/coul/dsf/omp}                 & half-Newton      & ---                  \\
10 & EXTRA-PAIR & \nolinkurl{born/coul/wolf/cs}             & \nolinkurl{born/coul/wolf/cs/omp}             & half-Newton      & ---                  \\
11 & EXTRA-PAIR & \nolinkurl{bpm/spring}                    & \nolinkurl{bpm/spring/omp}                    & half-Newton      & ---                  \\
12 & EXTRA-PAIR & \nolinkurl{dpd/fdt}                       & \nolinkurl{dpd/fdt/omp}                       & half-Newton      & random seed/thr      \\
13 & EXTRA-PAIR & \nolinkurl{dpd/fdt/energy}                & \nolinkurl{dpd/fdt/energy/omp}                & half-Newton      & random seed/thr      \\
14 & EXTRA-PAIR & \nolinkurl{lj/relres/soft}                & \nolinkurl{lj/relres/soft/omp}                & half-Newton      & ---                  \\
15 & EXTRA-PAIR & \nolinkurl{coul/diel}                     & \nolinkurl{coul/diel/omp}                     & half-Newton      & ---                  \\
16 & EXTRA-PAIR & \nolinkurl{momb}                          & \nolinkurl{momb/omp}                          & half-Newton      & ---                  \\
17 & EXTRA-PAIR & \nolinkurl{coul/streitz}                  & \nolinkurl{coul/streitz/omp}                  & half-Newton      & charge\_thr buf      \\
18 & EXTRA-PAIR & \nolinkurl{extep}                         & \nolinkurl{extep/omp}                         & half-Newton      & ---                  \\
19 & DIELECTRIC & \nolinkurl{lj/cut/coul/cut/dielectric}    & \nolinkurl{lj/cut/coul/cut/dielectric/omp}    & full (no Newton) & ef\_tally\_thr       \\
20 & DIELECTRIC & \nolinkurl{lj/cut/coul/long/dielectric}   & \nolinkurl{lj/cut/coul/long/dielectric/omp}   & full (no Newton) & ef\_tally\_thr       \\
21 & DIELECTRIC & \nolinkurl{lj/cut/coul/msm/dielectric}    & \nolinkurl{lj/cut/coul/msm/dielectric/omp}    & full (no Newton) & ef\_tally\_thr       \\
22 & DIELECTRIC & \nolinkurl{coul/long/dielectric}          & \nolinkurl{coul/long/dielectric/omp}          & full (no Newton) & ef\_tally\_thr       \\
23 & DIELECTRIC & \nolinkurl{lj/long/coul/long/dielectric}  & \nolinkurl{lj/long/coul/long/dielectric/omp}  & full (no Newton) & ef\_tally\_thr       \\
24 & DIELECTRIC & \nolinkurl{lj/cut/coul/debye/dielectric}  & \nolinkurl{lj/cut/coul/debye/dielectric/omp}  & full (no Newton) & ef\_tally\_thr       \\
25 & DIELECTRIC & \nolinkurl{coul/cut/dielectric}           & \nolinkurl{coul/cut/dielectric/omp}           & full (no Newton) & ef\_tally\_thr       \\
26 & KSPACE     & \nolinkurl{lj/long/coul/long}             & \nolinkurl{lj/long/coul/long/omp}             & half-Newton      & ---                  \\
27 & KSPACE     & \nolinkurl{buck/long/coul/long}           & \nolinkurl{buck/long/coul/long/omp}           & half-Newton      & ---                  \\
28 & KSPACE     & \nolinkurl{lj/long/tip4p/long}            & \nolinkurl{lj/long/tip4p/long/omp}            & half-Newton      & tip4p\_h\_thr        \\
29 & KSPACE     & \nolinkurl{lj/cut/tip4p/long}             & \nolinkurl{lj/cut/tip4p/long/omp}             & half-Newton      & tip4p\_h\_thr        \\
30 & KSPACE     & \nolinkurl{lj/long/dipole/long}           & \nolinkurl{lj/long/dipole/long/omp}           & half-Newton      & ---                  \\
31 & EXTRA-PAIR & \nolinkurl{buck/coul/cut/soft}            & \nolinkurl{buck/coul/cut/soft/omp}            & half-Newton      & ---                  \\
32 & EXTRA-PAIR & \nolinkurl{buck/coul/long/soft}           & \nolinkurl{buck/coul/long/soft/omp}           & half-Newton      & ---                  \\
33 & EXTRA-PAIR & \nolinkurl{lj/charmm/coul/long/soft}      & \nolinkurl{lj/charmm/coul/long/soft/omp}      & half-Newton      & ---                  \\
34 & EXTRA-PAIR & \nolinkurl{lj/charmm/coul/charmm/soft}    & \nolinkurl{lj/charmm/coul/charmm/soft/omp}    & half-Newton      & ---                  \\
35 & EXTRA-PAIR & \nolinkurl{coul/ctip}                     & \nolinkurl{coul/ctip/omp}                     & half-Newton      & charge\_iter\_thr    \\
\end{longtable}
\endgroup
\end{landscape}

\textbf{Notes}: - \emph{Loop type ``full''}: DIELECTRIC styles use a
full neighbor list (all neighbors i→j, not half-list) because the
dielectric boundary condition requires summing contributions from all
neighbors of i. No Newton's third law write-back is applied. -
\emph{ef\_tally\_thr}: The \texttt{ef\_tally\_thr()} function (added to
\texttt{src/OPENMP/thr\_omp.h}) accumulates per-thread electric field
contributions \texttt{efield{[}i{]}{[}k{]}} in addition to
energy/virial. - \emph{tip4p\_h\_thr}: For TIP4P water, the hydrogen
force redistribution (from the virtual M-site) is thread-safe only with
per-thread hydrogen force buffers. - \emph{random seed/thr}: DPD
stochastic forces use per-thread random number generators seeded with
\texttt{seed\ +\ thread\_id} to ensure deterministic but uncorrelated
thread-local sequences.

\begin{center}\rule{0.5\linewidth}{0.5pt}\end{center}

\subsection{Supplementary Table S2: A-3 Modification Log (Representative
Subset)}\label{supplementary-table-s2-a-3-modification-log-representative-subset}

Full log contains 102 entries. Representative entries spanning each
category:

% Table S2: scale-to-linewidth so file names with underscores fit
\begingroup
\setlength{\tabcolsep}{5pt}
\renewcommand{\arraystretch}{1.2}
\resizebox{\linewidth}{!}{%
\begin{tabular}{l l p{5.5cm} c p{2.6cm}}
\toprule
File & Category & Modification & Bit-id. & Pair improv. \\
\midrule
\nolinkurl{pair_lj_cut.cpp}               & Simple LJ      & \texttt{\_noalias} + \texttt{fxtmp/fytmp/fztmp}          & \checkmark & +4.2\% (4k FCC) \\
\nolinkurl{pair_lj_expand.cpp}            & Simple LJ      & same                                                      & \checkmark & +3.8\%          \\
\nolinkurl{pair_coul_cut.cpp}             & Simple Coul    & same                                                      & \checkmark & +4.1\%          \\
\nolinkurl{pair_coul_wolf.cpp}            & Comp.\ Coul    & same                                                      & \checkmark & +5.2\%          \\
\nolinkurl{pair_born_coul_dsf.cpp}        & Comp.\ Coul    & same                                                      & \checkmark & +5.7\%          \\
\nolinkurl{pair_born_coul_wolf_cs.cpp}    & Comp.\ Coul    & same                                                      & \checkmark & +4.9\%          \\
\nolinkurl{pair_lj_long_coul_long.cpp}    & Long-range     & same (real-space loop only)                               & \checkmark & +3.5\%          \\
\nolinkurl{pair_lj_cut_coul_long_dielectric.cpp} & DIELECTRIC & same + \texttt{ef\_tmp} for field accum          & \checkmark & +4.0\%          \\
\nolinkurl{pair_momb.cpp}                 & Specialty      & same                                                      & \checkmark & +7.5\%          \\
\nolinkurl{pair_coul_ctip.cpp}            & Charge-iter    & same (inner pair loop only)                               & \checkmark & +5.7\%          \\
\nolinkurl{pair_dispersion_d3.cpp}        & Multi-pass     & NOT applicable --- multi-pass loop structure              & n/a        & ---             \\
\nolinkurl{pair_eam.cpp}                  & Multi-body EAM & NOT applicable --- loop-carried $\phi_i$ accumulation     & n/a        & ---             \\
\nolinkurl{pair_tersoff.cpp}              & Multi-body     & NOT applicable --- three-body sums                        & n/a        & ---             \\
\nolinkurl{pair_snap.cpp}                 & ML potential   & NOT applicable --- internal vectorization                 & n/a        & ---             \\
\bottomrule
\end{tabular}
}% end resizebox
\endgroup

\textbf{Note on non-applicable styles}: Multi-body potentials (EAM,
Tersoff, SW, ReaxFF) use multi-pass loops where force contributions from
one atom depend on previously visited neighbors; the \texttt{fxtmp}
pattern requires a single-pass inner loop. ML potentials (SNAP, ACE,
DeePMD) implement their own BLAS-based or SIMD-specialized inner
operations that supersede simple pointer qualification.

\begin{center}\rule{0.5\linewidth}{0.5pt}\end{center}

\subsection{Supplementary Table S3: OPLS-AA Force Field Parameters for
Battery Electrolyte
Solvents}\label{supplementary-table-s3-opls-aa-force-field-parameters-for-battery-electrolyte-solvents}

All parameters follow Jorgensen \emph{et al.} (1996) OPLS-AA
conventions. LJ potential: U\_LJ = 4ε{[}(σ/r)¹²−(σ/r)⁶{]}; harmonic
bond: U\_bond = K(r−r₀)²; harmonic angle: U\_angle = K(θ−θ₀)²; OPLS
dihedral: U\_torsion = ½V₁(1+cos φ) + ½V₂(1−cos 2φ) + ½V₃(1+cos 3φ) +
½V₄(1−cos 4φ).

\subsubsection{Atom Types (shared across EC and
PC)}\label{atom-types-shared-across-ec-and-pc}

{\def\LTcaptype{none} % do not increment counter
\begin{longtable}[]{@{}
  >{\raggedright\arraybackslash}p{(\linewidth - 8\tabcolsep) * \real{0.1618}}
  >{\raggedright\arraybackslash}p{(\linewidth - 8\tabcolsep) * \real{0.1912}}
  >{\raggedright\arraybackslash}p{(\linewidth - 8\tabcolsep) * \real{0.1176}}
  >{\raggedright\arraybackslash}p{(\linewidth - 8\tabcolsep) * \real{0.2059}}
  >{\raggedright\arraybackslash}p{(\linewidth - 8\tabcolsep) * \real{0.3235}}@{}}
\toprule\noalign{}
\begin{minipage}[b]{\linewidth}\raggedright
Atom type
\end{minipage} & \begin{minipage}[b]{\linewidth}\raggedright
Description
\end{minipage} & \begin{minipage}[b]{\linewidth}\raggedright
σ (Å)
\end{minipage} & \begin{minipage}[b]{\linewidth}\raggedright
ε (kcal/mol)
\end{minipage} & \begin{minipage}[b]{\linewidth}\raggedright
Partial charge q (e)
\end{minipage} \\
\midrule\noalign{}
\endhead
\bottomrule\noalign{}
\endlastfoot
C\_co3 & Carbonyl C in carbonate & 3.750 & 0.1050 & +0.76 \\
O\_co3 & Carbonyl O (C=O) & 2.960 & 0.2100 & −0.50 \\
O\_ete & Ether O (ring C−O−C) & 3.000 & 0.1700 & −0.35 \\
C\_sp3 & Aliphatic CH₂ in ring & 3.500 & 0.0660 & +0.10 \\
H\_c & H on aliphatic C & 2.500 & 0.0300 & +0.06 \\
\end{longtable}
}

\textbf{EC molecule charge neutrality}: C\_co3 + 2×O\_co3 + 2×O\_ete +
2×C\_sp3 + 4×H\_c = +0.76 + 2×(−0.50) + 2×(−0.35) + 2×(+0.10) +
4×(+0.06) = +0.76 − 1.00 − 0.70 + 0.20 + 0.24 = 0.00 (neutral)

\subsubsection{Atom Types (DME:
1,2-dimethoxyethane)}\label{atom-types-dme-12-dimethoxyethane}

{\def\LTcaptype{none} % do not increment counter
\begin{longtable}[]{@{}
  >{\raggedright\arraybackslash}p{(\linewidth - 8\tabcolsep) * \real{0.1618}}
  >{\raggedright\arraybackslash}p{(\linewidth - 8\tabcolsep) * \real{0.1912}}
  >{\raggedright\arraybackslash}p{(\linewidth - 8\tabcolsep) * \real{0.1176}}
  >{\raggedright\arraybackslash}p{(\linewidth - 8\tabcolsep) * \real{0.2059}}
  >{\raggedright\arraybackslash}p{(\linewidth - 8\tabcolsep) * \real{0.3235}}@{}}
\toprule\noalign{}
\begin{minipage}[b]{\linewidth}\raggedright
Atom type
\end{minipage} & \begin{minipage}[b]{\linewidth}\raggedright
Description
\end{minipage} & \begin{minipage}[b]{\linewidth}\raggedright
σ (Å)
\end{minipage} & \begin{minipage}[b]{\linewidth}\raggedright
ε (kcal/mol)
\end{minipage} & \begin{minipage}[b]{\linewidth}\raggedright
Partial charge q (e)
\end{minipage} \\
\midrule\noalign{}
\endhead
\bottomrule\noalign{}
\endlastfoot
C\_sp3(CH₃) & Terminal methyl C & 3.500 & 0.0660 & +0.022 \\
O\_ete & Ether O & 3.000 & 0.1700 & −0.400 \\
C\_sp3(CH₂) & Central methylene C & 3.500 & 0.0660 & +0.164 \\
H\_c(CH₃) & H on methyl & 2.500 & 0.0300 & +0.034 \\
H\_c(CH₂) & H on methylene & 2.500 & 0.0300 & +0.056 \\
\end{longtable}
}

\textbf{DME molecule charge neutrality}: 2×C\_sp3(CH₃) + 2×O\_ete +
2×C\_sp3(CH₂) + 6×H\_c(CH₃) + 4×H\_c(CH₂) = 2×(+0.022) + 2×(−0.400) +
2×(+0.164) + 6×(+0.034) + 4×(+0.056) = 0.044 − 0.800 + 0.328 + 0.204 +
0.224 = 0.000 (neutral)

\subsubsection{Bond Parameters}\label{bond-parameters}

{\def\LTcaptype{none} % do not increment counter
\begin{longtable}[]{@{}lll@{}}
\toprule\noalign{}
Bond & K (kcal/mol/Ų) & r₀ (Å) \\
\midrule\noalign{}
\endhead
\bottomrule\noalign{}
\endlastfoot
C\_co3−O\_co3 & 570.0 & 1.229 \\
C\_co3−O\_ete & 415.0 & 1.344 \\
C\_sp3−O\_ete & 320.0 & 1.430 \\
C\_sp3−C\_sp3 & 224.0 & 1.529 \\
C\_sp3−H\_c & 340.0 & 1.090 \\
\end{longtable}
}

\subsubsection{Angle Parameters}\label{angle-parameters}

{\def\LTcaptype{none} % do not increment counter
\begin{longtable}[]{@{}lll@{}}
\toprule\noalign{}
Angle & K (kcal/mol/rad²) & θ₀ (°) \\
\midrule\noalign{}
\endhead
\bottomrule\noalign{}
\endlastfoot
O\_co3−C\_co3−O\_ete & 90.0 & 124.4 \\
C\_co3−O\_ete−C\_sp3 & 70.0 & 109.5 \\
O\_ete−C\_sp3−C\_sp3 & 70.0 & 111.5 \\
O\_ete−C\_sp3−H\_c & 55.0 & 109.5 \\
H\_c−C\_sp3−H\_c & 35.0 & 107.8 \\
C\_sp3−C\_sp3−H\_c & 37.5 & 110.7 \\
\end{longtable}
}

\subsubsection{Dihedral Parameters (OPLS style, V₁−V₄ in
kcal/mol)}\label{dihedral-parameters-opls-style-vux2081vux2084-in-kcalmol}

{\def\LTcaptype{none} % do not increment counter
\begin{longtable}[]{@{}lllll@{}}
\toprule\noalign{}
Dihedral & V₁ & V₂ & V₃ & V₄ \\
\midrule\noalign{}
\endhead
\bottomrule\noalign{}
\endlastfoot
O\_co3−C\_co3−O\_ete−C\_sp3 & 1.950 & 2.250 & 0.000 & 0.000 \\
C\_co3−O\_ete−C\_sp3−C\_sp3 & 0.650 & −0.250 & 0.670 & 0.000 \\
C\_co3−O\_ete−C\_sp3−H\_c & 0.000 & 0.000 & 0.760 & 0.000 \\
O\_ete−C\_sp3−C\_sp3−O\_ete & −0.550 & −0.550 & 0.000 & 0.000 \\
O\_ete−C\_sp3−C\_sp3−H\_c & 0.000 & 0.000 & 0.468 & 0.000 \\
H\_c−C\_sp3−C\_sp3−H\_c & 0.000 & 0.000 & 0.318 & 0.000 \\
\end{longtable}
}

\begin{center}\rule{0.5\linewidth}{0.5pt}\end{center}

\subsection{Supplementary Table S4: Li⁺ Parameters and
Lorentz--Berthelot Mixed Pair
Coefficients}\label{supplementary-table-s4-li-parameters-and-lorentzberthelot-mixed-pair-coefficients}

\textbf{Li⁺ source}: Joung \& Cheatham (2008), SPC/E water
parameterization.\\
\textbf{Mixing rules}: σ₁₂ = (σ₁ + σ₂)/2 (arithmetic); ε₁₂ = √(ε₁ × ε₂)
(geometric).

{\def\LTcaptype{none} % do not increment counter
\begin{longtable}[]{@{}lll@{}}
\toprule\noalign{}
Parameter & Li⁺ & Reference \\
\midrule\noalign{}
\endhead
\bottomrule\noalign{}
\endlastfoot
σ (Å) & 1.4094 & Joung \& Cheatham 2008 \\
ε (kcal/mol) & 0.3367 & Joung \& Cheatham 2008 \\
q (e) & +1.0 & formal charge \\
\end{longtable}
}

\subsubsection{Lorentz--Berthelot Mixed Coefficients: Li⁺ × OPLS-AA
Solvent
Atoms}\label{lorentzberthelot-mixed-coefficients-li-opls-aa-solvent-atoms}

{\def\LTcaptype{none} % do not increment counter
\begin{longtable}[]{@{}llll@{}}
\toprule\noalign{}
Pair & σ\_mix (Å) & ε\_mix (kcal/mol) & Context \\
\midrule\noalign{}
\endhead
\bottomrule\noalign{}
\endlastfoot
Li--Li & 1.4094 & 0.33670 & self-pair (irrelevant, single Li+) \\
Li--C\_co3 & 2.5797 & 0.18802 & EC/PC carbonyl carbon \\
Li--O\_co3 & 2.1847 & 0.26591 & EC/PC carbonyl oxygen \\
Li--O\_ete & 2.2047 & 0.23921 & EC/PC/DME ether oxygen \\
Li--C\_sp3 & 2.4547 & 0.14898 & EC/PC/DME aliphatic carbon \\
Li--H\_c & 1.9547 & 0.10049 & all H atoms \\
\end{longtable}
}

\textbf{Note on Coulombic interactions}: The
\texttt{lj/cut/coul/long/soft} potential scales \emph{both} the LJ and
Coulomb interactions by λ simultaneously. The Li⁺ charge (+1.0 e)
interacts with OPLS-AA partial charges via PPPM long-range Coulomb. The
dominant contribution to dU/dλ is the Coulombic Li⁺--O\_co3 and
Li⁺--O\_ete attraction (solvation coordination shell). A +1 net charge
in the simulation box is handled by PPPM with a uniform background
charge correction, which introduces a small but systematic correction to
the absolute ΔG (estimated \textless{} 0.1 kcal/mol for the box sizes
used) without affecting relative comparisons between solvents {[}ref:
Hummer et al.~1996, J. Phys. Chem.{]}.

\begin{center}\rule{0.5\linewidth}{0.5pt}\end{center}

\subsection{Supplementary Table S5: OPLS-AA FEP Results --- dU/dλ by
Lambda
Window}\label{supplementary-table-s5-opls-aa-fep-results-dudux3bb-by-lambda-window}

Values shown are mean ± SEM (kcal/mol) over 40 ps production (20,000
steps × 2 fs/step) from individual-window equilibration (50 ps NPT + 50
ps NVT for EC/PC; NVT-only for DME). All three estimators (TI, BAR,
MBAR) applied; TI uses the trapezoidal rule over 21 points.

\subsubsection{EC (ethylene carbonate) +
Li⁺}\label{ec-ethylene-carbonate-li}

Production: 40 ps (20,000 × 2 fs steps), 101 samples/window. Protocol:
NPT 100 ps + NVT 25 ps equilibration per window.

{\def\LTcaptype{none} % do not increment counter
\begin{longtable}[]{@{}llll@{}}
\toprule\noalign{}
λ & ⟨dU/dλ⟩ (kcal/mol) & SEM & n\_frames \\
\midrule\noalign{}
\endhead
\bottomrule\noalign{}
\endlastfoot
0.00 & −0.29 & 0.03 & 101 \\
0.05 & −2.88 & 0.27 & 101 \\
0.10 & −8.99 & 0.90 & 101 \\
0.15 & −56.61 & 2.19 & 101 \\
0.20 & −109.42 & 1.48 & 101 \\
0.25 & −161.04 & 1.66 & 101 \\
0.30 & −230.18 & 1.06 & 101 \\
0.35 & −292.22 & 1.26 & 101 \\
0.40 & −378.43 & 1.30 & 101 \\
0.45 & −454.16 & 1.71 & 101 \\
0.50 & −530.22 & 1.41 & 101 \\
0.55 & −80.04 & 0.99 & 101 \\
0.60 & −89.00 & 1.10 & 101 \\
0.65 & −106.08 & 1.04 & 101 \\
0.70 & −122.67 & 0.95 & 101 \\
0.75 & −139.93 & 0.97 & 101 \\
0.80 & −161.66 & 0.94 & 101 \\
0.85 & −179.61 & 1.21 & 101 \\
0.90 & −207.37 & 1.32 & 101 \\
0.95 & −249.79 & 1.01 & 101 \\
1.00 & −267.21 & 1.12 & 101 \\
\end{longtable}
}

\textbf{Note on non-monotone integrand}: The dU/dλ curve exhibits a
characteristic peak at λ=0.50 (−530 kcal/mol) followed by a sharp drop
to −80 kcal/mol at λ=0.55. This arises from the simultaneous LJ
soft-core (n=2, α\_LJ=0.5) and Coulomb soft-core (α\_C=0.3) coupling in
\texttt{lj/cut/coul/long/soft}: the LJ soft-core contributes maximally
near λ≈0.5, while the Coulomb term dominates at λ\textgreater0.5. The
abrupt transition at λ=0.50→0.55 is the primary driver of the TI--BAR
discrepancy (15.6 kcal/mol, 9\%) and motivates the extended 200
ps/window run.

\subsubsection{PC (propylene carbonate) +
Li⁺}\label{pc-propylene-carbonate-li}

Production: 40 ps (20,000 × 2 fs steps), 101 samples/window. Protocol:
NPT 100 ps + NVT 25 ps equilibration per window.

{\def\LTcaptype{none} % do not increment counter
\begin{longtable}[]{@{}llll@{}}
\toprule\noalign{}
λ & ⟨dU/dλ⟩ (kcal/mol) & SEM & n\_frames \\
\midrule\noalign{}
\endhead
\bottomrule\noalign{}
\endlastfoot
0.00 & −0.26 & 0.03 & 101 \\
0.05 & −3.33 & 0.32 & 101 \\
0.10 & −8.43 & 0.88 & 101 \\
0.15 & −36.81 & 2.55 & 101 \\
0.20 & −106.69 & 1.65 & 101 \\
0.25 & −166.00 & 1.16 & 101 \\
0.30 & −218.07 & 1.51 & 101 \\
0.35 & −295.32 & 1.37 & 101 \\
0.40 & −378.32 & 1.23 & 101 \\
0.45 & −463.10 & 0.99 & 101 \\
0.50 & −519.57 & 1.63 & 101 \\
0.55 & −75.23 & 1.10 & 101 \\
0.60 & −88.81 & 1.05 & 101 \\
0.65 & −111.05 & 1.05 & 101 \\
0.70 & −122.34 & 0.97 & 101 \\
0.75 & −139.85 & 1.07 & 101 \\
0.80 & −154.54 & 1.01 & 101 \\
0.85 & −177.73 & 1.09 & 101 \\
0.90 & −209.25 & 1.40 & 101 \\
0.95 & −244.26 & 1.40 & 101 \\
1.00 & −276.25 & 1.05 & 101 \\
\end{longtable}
}

\subsubsection{DME (1,2-dimethoxyethane) +
Li⁺}\label{dme-12-dimethoxyethane-li}

Production: 40 ps (20,000 × 2 fs steps), 101 samples/window. Protocol:
NVT-only (staged: nve/limit → NVT ramp → NVT 35 ps equilibration per
window).

{\def\LTcaptype{none} % do not increment counter
\begin{longtable}[]{@{}llll@{}}
\toprule\noalign{}
λ & ⟨dU/dλ⟩ (kcal/mol) & SEM & n\_frames \\
\midrule\noalign{}
\endhead
\bottomrule\noalign{}
\endlastfoot
0.00 & −0.26 & 0.03 & 101 \\
0.05 & −2.54 & 0.30 & 101 \\
0.10 & −8.13 & 0.92 & 101 \\
0.15 & −32.15 & 2.62 & 101 \\
0.20 & −92.66 & 1.70 & 101 \\
0.25 & −141.94 & 1.36 & 101 \\
0.30 & −192.26 & 1.23 & 101 \\
0.35 & −248.70 & 1.33 & 101 \\
0.40 & −314.17 & 1.11 & 101 \\
0.45 & −375.96 & 1.23 & 101 \\
0.50 & −424.57 & 1.35 & 101 \\
0.55 & −73.05 & 0.95 & 101 \\
0.60 & −85.54 & 1.13 & 101 \\
0.65 & −95.54 & 0.82 & 101 \\
0.70 & −109.23 & 0.90 & 101 \\
0.75 & −124.77 & 1.28 & 101 \\
0.80 & −159.45 & 1.13 & 101 \\
0.85 & −180.73 & 1.09 & 101 \\
0.90 & −199.54 & 0.97 & 101 \\
0.95 & −217.50 & 0.91 & 101 \\
1.00 & −237.92 & 0.83 & 101 \\
\end{longtable}
}

\textbf{Note on DME dU/dλ profile}: Peak at λ=0.50 (−424.6 kcal/mol) is
smaller than EC/PC (\textasciitilde−520--530 kcal/mol), reflecting the
weaker OPLS-AA oxygen partial charges on ether oxygens (−0.40 e)
vs.~carbonate carbonyl oxygens (−0.50 e). The bidentate coordination of
DME (two oxygens/molecule) partially compensates for the weaker
per-oxygen interaction.

\subsubsection{Free energy summary (40 ps/window short
run)}\label{free-energy-summary-40-pswindow-short-run}

\emph{These values correspond to the 40 ps/window production run. The
200 ps/window extended-run values (Table 7 in main text) are: EC −183.7
± 0.1 (TI), −168.4 (BAR); PC −182.1 ± 0.1 (TI), −167.2 (BAR); DME −159.4
± 0.1 (TI), −147.0 (BAR).}

{\def\LTcaptype{none} % do not increment counter
\begin{longtable}[]{@{}
  >{\raggedright\arraybackslash}p{(\linewidth - 8\tabcolsep) * \real{0.1139}}
  >{\raggedright\arraybackslash}p{(\linewidth - 8\tabcolsep) * \real{0.3038}}
  >{\raggedright\arraybackslash}p{(\linewidth - 8\tabcolsep) * \real{0.3165}}
  >{\raggedright\arraybackslash}p{(\linewidth - 8\tabcolsep) * \real{0.1519}}
  >{\raggedright\arraybackslash}p{(\linewidth - 8\tabcolsep) * \real{0.1139}}@{}}
\toprule\noalign{}
\begin{minipage}[b]{\linewidth}\raggedright
Solvent
\end{minipage} & \begin{minipage}[b]{\linewidth}\raggedright
ΔG\_TI 40 ps (kcal/mol)
\end{minipage} & \begin{minipage}[b]{\linewidth}\raggedright
ΔG\_BAR 40 ps (kcal/mol)
\end{minipage} & \begin{minipage}[b]{\linewidth}\raggedright
TI--BAR gap
\end{minipage} & \begin{minipage}[b]{\linewidth}\raggedright
Windows
\end{minipage} \\
\midrule\noalign{}
\endhead
\bottomrule\noalign{}
\endlastfoot
EC & \textbf{−184.7 ± 0.3} & \textbf{−169.1} & 15.6 (9\%) & 21/21 \\
PC & \textbf{−182.8 ± 0.3} & \textbf{−167.8} & 15.0 (8\%) & 21/21 \\
DME & \textbf{−159.9 ± 0.3} & \textbf{−147.4} & 12.5 (8\%) & 21/21 \\
CG (reference) & −0.257 ± 0.004 & −0.265 & 0.008 (3\%) & 21/21 \\
\end{longtable}
}

Relative solvation strengths (TI): ΔΔΔG(EC−DME) = 24.8 kcal/mol;
ΔΔΔG(PC−DME) = 22.9 kcal/mol; ΔΔΔG(EC−PC) = 1.9 kcal/mol.

\emph{Physical interpretation}: ΔG is the free energy of coupling Li⁺ to
the solvent (λ: 0→1). Solvation free energy = −ΔG\_coupling. More
negative ΔG\_coupling indicates stronger solvation. Expected ordering:
EC \textgreater{} PC \textgreater{} DME based on dielectric constant
(ε\_r: EC=89.8, PC=64.9, DME=7.2) and Li⁺ coordination in quantum
chemistry studies.

\begin{center}\rule{0.5\linewidth}{0.5pt}\end{center}

\subsection{Supplementary Methods S1: Claude Code Prompt
Templates}\label{supplementary-methods-s1-claude-code-prompt-templates}

\subsubsection{S1.1 Template for Half-List OMP Port (EXTRA-PAIR
styles)}\label{s1.1-template-for-half-list-omp-port-extra-pair-styles}

\begin{verbatim}
Context:
  Working directory: lammps-src/src/
  Task: Add OMP variant for pair style [STYLE_NAME]
  Source files: EXTRA-PAIR/pair_[style_file].{h,cpp}
  Reference implementation: OPENMP/pair_lj_cut_omp.{h,cpp}
  Reference source: pair_lj_cut.{h,cpp}

Rules for half-list Newton-on style:
  1. File pair_[style_file]_omp.h:
     - class Pair[StyleClass]OMP : public Pair[StyleClass], public ThrOMP
     - Constructor sets "no_virial_fdotr_compute = 1"
     - Override: compute(int, int), settings(int, char**), init_one(int, int)
     - Keep all data members from parent; no new members except inherited ThrOMP

  2. File pair_[style_file]_omp.cpp:
     - compute() parallel region pattern (copy from pair_lj_cut_omp.cpp:80-155):
       a. OpenMP parallel over outer atom loop
       b. thr->get_f() provides per-thread force buffer
       c. Dispatch: eval<0,0,1>(), eval<0,1,1>(), eval<1,0,1>(), eval<1,1,1>()
          (EVFLAG×EFLAG×NEWTON_PAIR combinations)
       d. reduce_thr(this, eflag, vflag, thr) at end of parallel region
     - eval<EVFLAG,EFLAG,NEWTON_PAIR>() method:
       a. Reads: atom->x, atom->type, list->ilist etc. from parent class
       b. Inner neighbor loop uses list->firstneigh[ii] etc.
       c. Force accumulation into fxtmp/fytmp/fztmp (not f[i][k] directly)
       d. Energy/virial: ev_tally_thr(this, i, j, nlocal, NEWTON_PAIR,
                                      evdwl, ecoul, fpair, delx, dely, delz, thr)
       e. Newton write: if NEWTON_PAIR: fj[0]-=fpairx, fj[1]-=fpairy, fj[2]-=fpairz
  
  3. CMakeLists.txt (src/OPENMP/CMakeLists.txt):
     - Add pair_[style_file]_omp.{h,cpp} in alphabetical order
  
  4. src/OPENMP/style_pair.h or the auto-generated pair styles list:
     - Add:  PairStyle(lj/cut/[name]/omp, Pair[Class]OMP)

Verification:
  Compile with: cmake --build build --config Release --target lmp
  Run serial: lmp.exe -in tools/tests/in.[style]_test -log none > serial.out
  Run OMP:    lmp.exe -sf omp -pk omp 4 -in tools/tests/in.[style]_test -log none > omp.out
  Verify: python tools/tests/compare_thermo.py serial.out omp.out
          (must print "PASS: bit-identical" for all columns)
\end{verbatim}

\subsubsection{S1.2 Template for DIELECTRIC Full-List OMP
Port}\label{s1.2-template-for-dielectric-full-list-omp-port}

\begin{verbatim}
Context:
  Task: Add OMP variant for DIELECTRIC pair style [STYLE_NAME]
  Source: DIELECTRIC/pair_[style_file].{h,cpp}
  Reference OMP structure: OPENMP/pair_lj_cut_coul_long_omp.cpp
  Reference DIELECTRIC physics: DIELECTRIC/pair_lj_cut_coul_cut_dielectric.cpp

Special rules for DIELECTRIC full-list:
  1. No Newton's third law (full list, each pair counted once from i's side only)
     → remove all "if (NEWTON_PAIR) {fj[k] -= ...}" blocks
     → template has only one NEWTON_PAIR=0 specialization
  2. Per-atom dielectric epsilon: read atom->epsilon array
     epsj = (epsiloni + epsilon[jtype]) * 0.5
  3. Electric field accumulation:
     ef_tally_thr(this, i, j, nlocal, evdwl_ij, ex, ey, ez, thr)
     where (ex,ey,ez) is the field contribution from pair i-j
  4. Thread-safe efield accumulation uses per-thread ef_buf in ThrOMP
  5. Outer atom loop is over ALL local atoms (not ghost-list subset)
     neigh_list->ilist covers i in [0, nlocal)

Verification: same as S1.1 but also check efield[i][k] values match serial.
\end{verbatim}

\subsubsection{S1.3 Template for A-3 Restrict/fxtmp
Backport}\label{s1.3-template-for-a-3-restrictfxtmp-backport}

\begin{verbatim}
Context:
  Task: Apply A-3 restrict/fxtmp optimization to src/pair_[name].cpp
  Target function: compute(int eflag, int vflag)
  Reference before-state: src/pair_lj_cut.cpp (lines 69-139)
  Reference after-state:  src/OPENMP/pair_lj_cut_omp.cpp (lines 99-140, eval<>)

Transformation rules (apply ONLY to the inner Newton-pair loop; do NOT touch
multi-body loops, charge equilibration loops, or indirect-index patterns):

  1. Immediately before the outer loop (for i = 0; i < inum; i++):
     Add:
       double * _noalias const f0 = atom->f[0];
     Replace existing  double **f = atom->f;  (keep x, v, other **ptrs unchanged)

  2. At the START of the outer loop body (inside "for ii = 0..."):
     Add:
       const int i   = ilist[ii];
       double fxtmp  = 0.0;
       double fytmp  = 0.0;
       double fztmp  = 0.0;
       double * fi   = f0 + 3*i;   // pointer arithmetic for AoS layout

  3. In the INNER loop: replace
       f[i][0] += delx*fpair;
       f[i][1] += dely*fpair;
       f[i][2] += delz*fpair;
     with:
       fxtmp += delx*fpair;
       fytmp += dely*fpair;
       fztmp += delz*fpair;

  4. Keep Newton write-back (fj[k] -= fk) UNCHANGED.

  5. After the inner loop (jj loop), before next outer iteration:
     Add:
       fi[0] += fxtmp;
       fi[1] += fytmp;
       fi[2] += fztmp;

CRITICAL constraint: The floating-point operation ORDER for each (delx*fpair),
(dely*fpair), (delz*fpair) must be unchanged relative to the original code.
Only the partial accumulation into f[i][k] is deferred to after the inner loop.

Verification:
  Run: lmp.exe -in tools/tests/in.[style]_test > before.out  (original binary)
       lmp.exe -in tools/tests/in.[style]_test > after.out   (modified binary)
  Check: python tools/tests/compare_thermo.py before.out after.out
         Must print "PASS: bit-identical"
\end{verbatim}

\begin{center}\rule{0.5\linewidth}{0.5pt}\end{center}

\subsection{Supplementary Methods S2: OPLS-AA FEP Case Study --- Full
Protocol}\label{supplementary-methods-s2-opls-aa-fep-case-study-full-protocol}

\subsubsection{S2.1 System Construction}\label{s2.1-system-construction}

All-atom data files were generated by a custom Python script
(\texttt{tools/fep/data/build\_opls\_systems.py}) without external tools
(no Packmol required). The packing algorithm:

\begin{enumerate}
\def\labelenumi{\arabic{enumi}.}
\tightlist
\item
  Computes target box length from molecular density: L = (N × M\_mol /
  (N\_A × ρ\_target))\^{}(1/3)
\item
  Generates a 3D integer grid of N\_mol sites within a box of L\_scale ×
  L
\item
  Assigns each molecule a random SO(3) rotation via Rodrigues formula
\item
  Adds a uniform random jitter ±0.1 L\_scale × L / N\_sites\^{}(1/3) per
  site
\item
  Places atoms relative to a pre-defined geometry template with bond
  lengths at equilibrium values
\end{enumerate}

Molecule geometries: - \textbf{EC}: Regular pentagon ring
(C₁--O₁--C₂--O₂--C₃, circumradius R=1.216 Å), carbonyl O at R=1.206 Å
from C₁, H atoms at ±z=1.09 Å from CH₂ carbons - \textbf{PC}: EC ring
with C₃ replaced by a CH group bearing a −CH₃ substituent at (0.715,
−1.609, −1.408) Å relative to the ring center - \textbf{DME}: All-trans
backbone; C₁H₃−O₁−C₂H₂−C₃H₂−O₂−C₄H₃ along the x-axis; methyl H atoms
positioned via tetrahedral geometry

The \texttt{lj/cut/coul/long/soft} pair style was applied to ALL Li⁺
cross-pairs with the lambda parameter; solvent--solvent interactions
used standard \texttt{lj/cut/coul/long} (not soft-core, no lambda
coupling). This is the standard single-topology FEP approach for
charging/annihilating a single ion.

\subsubsection{S2.2 Equilibration
Protocol}\label{s2.2-equilibration-protocol}

\textbf{EC and PC} (601 atoms / 651 atoms): 1. Energy minimization:
conjugate-gradient, etol=10⁻⁴, ftol=10⁻⁶, max 5000 iterations / 50000
force evaluations 2. Velocity initialization: 300 K, Gaussian
distribution, seed=12345 3. NPT equilibration: 50,000 steps (100 ps),
Nosé--Hoover T=300 K (τ=200 fs), iso P=1 atm (τ=1000 fs) 4. NVT
equilibration: 25,000 steps (50 ps), Nosé--Hoover T=300 K (τ=200 fs)

\textbf{DME} (961 atoms; NVT-only to avoid PPPM out-of-range errors at
the target density): 1. Energy minimization: etol=10⁻⁵, ftol=10⁻⁷, max
10000 iterations / 100000 force evaluations 2. Velocity init: 100 K
(gentle start) 3. NVE/limit: 5000 steps with max displacement 0.05
Å/step + temp/rescale 100→200 K 4. NVT ramp: 10,000 steps (15 ps),
200→300 K (τ=200 fs), timestep=1.5 fs 5. NVT production equil: 35,000
steps (70 ps), 300 K, timestep=2 fs

\subsubsection{S2.3 FEP Production
Protocol}\label{s2.3-fep-production-protocol}

For each of the 21 lambda windows (λ = 0.00, 0.05, \ldots, 1.00) and
each of the 3 solvents: - Starting configuration: re-generated from
scratch for each window (independent equilibration) - Post-equilibration
NVT: 5000 steps (10 ps) - Production: 20,000 steps (40 ps), thermo
output every 200 steps (100 frames/window) - PPPM: accuracy 10⁻⁴, mesh
automatically determined - FEP compute:
\texttt{compute\ FEP\ all\ fep\ 300.0\ pair\ lj/cut/coul/long/soft\ lambda\ 1\ *\ 0.01}
- Output: step, λ, dU/dλ, exp(−ΔU/kT) via
\texttt{thermo\_style\ custom\ step\ v\_LAMv\ v\_dUdL\ c\_FEP{[}2{]}}

\subsubsection{S2.4 Free Energy
Analysis}\label{s2.4-free-energy-analysis}

TI: trapezoidal rule over 21 λ points\\
BAR: adjacent-window estimator using
\texttt{pymbar.other\_estimators.bar()} (pymbar 4.x)\\
MBAR: full multi-state estimator using linear potential energy
approximation

The TI error is propagated from the per-window standard error of the
mean (SEM), combined with the trapezoidal quadrature weights as
described in Shirts \& Chodera (2008).

\subsubsection{S2.5 Charge Neutrality and
PPPM}\label{s2.5-charge-neutrality-and-pppm}

The simulation box contains +1 net charge (from Li⁺). PPPM in LAMMPS
handles this by adding a compensating uniform background charge density
(the ``jellium'' background), which introduces a spurious
self-interaction energy E\_corr = −q²/(2ε₀V) × (dimensionful prefactor).
For a box of V ≈ 30³ ų and q = +1e, this correction is approximately
0.02--0.05 kcal/mol and affects all λ windows equally, canceling in the
ΔG\_coupling calculation. For absolute solvation free energies, a
finite-size correction following Hummer et al.~(1996) should be applied;
this is outside the scope of the present method demonstration.

\begin{center}\rule{0.5\linewidth}{0.5pt}\end{center}

\subsection{Supplementary Figure S1: Full OMP Scaling Benchmarks for New
Styles}\label{supplementary-figure-s1-full-omp-scaling-benchmarks-for-new-styles}

\emph{{[}Figure generated by the Run-Benchmark.ps1 harness; stored as
tools/plots/figure\_s1\_omp\_scaling.pdf{]}}

OMP scaling benchmark for three representative newly ported pair styles:
\texttt{born/coul/dsf}, \texttt{lj/class2/soft}, and
\texttt{nm/cut/split}. Each style was benchmarked on an 864-atom FCC
lattice (500 NVE steps, N=3 runs) at 1, 2, 4, and 8 OMP threads using
\texttt{-sf\ omp\ -pk\ omp\ N}. Data correspond to Table 4 in the main
text (same 500-step runs).

\textbf{Panel (a)}: Loop time (s) vs.~OMP thread count for all three
styles. Error bars show ±1 standard deviation over 3 runs.

\textbf{Panel (b)}: Speedup (serial loop time / N-thread loop time)
vs.~OMP thread count. Dashed line shows ideal linear scaling. Amdahl-law
fits (parallelizable fraction f) overlay the measured data,
demonstrating that pair-dominated styles (pair fraction \textgreater{}
90\%) approach but do not reach the ideal limit due to fixed overhead in
neighbor building and output.

\textbf{Panel (c)}: Pair time fraction (pair time / loop time)
vs.~thread count, confirming that the non-parallelized overhead grows as
a relative fraction with increasing thread count---the primary cause of
sub-linear scaling at 8 threads.

Measured speedup values (500-step runs, N=3, mean ± 1σ). These values
are the same dataset as Table 4 in the main text.

{\def\LTcaptype{none} % do not increment counter
\begin{longtable}[]{@{}
  >{\raggedright\arraybackslash}p{(\linewidth - 8\tabcolsep) * \real{0.2000}}
  >{\raggedright\arraybackslash}p{(\linewidth - 8\tabcolsep) * \real{0.2000}}
  >{\raggedright\arraybackslash}p{(\linewidth - 8\tabcolsep) * \real{0.2000}}
  >{\raggedright\arraybackslash}p{(\linewidth - 8\tabcolsep) * \real{0.2000}}
  >{\raggedright\arraybackslash}p{(\linewidth - 8\tabcolsep) * \real{0.2000}}@{}}
\toprule\noalign{}
\begin{minipage}[b]{\linewidth}\raggedright
Style
\end{minipage} & \begin{minipage}[b]{\linewidth}\raggedright
4t speedup
\end{minipage} & \begin{minipage}[b]{\linewidth}\raggedright
8t speedup
\end{minipage} & \begin{minipage}[b]{\linewidth}\raggedright
Pair fraction (1t)
\end{minipage} & \begin{minipage}[b]{\linewidth}\raggedright
Serial fraction \emph{f}
\end{minipage} \\
\midrule\noalign{}
\endhead
\bottomrule\noalign{}
\endlastfoot
\texttt{born/coul/dsf} & \textbf{2.92×} & \textbf{3.96×} & 92.9\% &
12\% \\
\texttt{lj/class2/soft} & \textbf{3.14×} & \textbf{3.75×} & 94.8\% &
9\% \\
\texttt{nm/cut/split} & \textbf{2.63×} & \textbf{2.81×} & 85.9\% &
17\% \\
\end{longtable}
}

The serial fraction \emph{f} was estimated from Amdahl's law fitted to
the measured 4-thread speedup: f = (4/S₄ − 1)/3. \texttt{nm/cut/split}
has the largest \emph{f} (17\%) because approximately 10\% of loop time
is spent in neighbor-list building (\texttt{Neigh} section), which does
not benefit from OMP pair parallelism.

\begin{center}\rule{0.5\linewidth}{0.5pt}\end{center}

\subsection{Supplementary Figure S2: NVE Energy Conservation
Comparison}\label{supplementary-figure-s2-nve-energy-conservation-comparison}

\emph{{[}Figure description{]}}

For verification of all new OMP variants, 250-step NVE simulations were
run with \texttt{thermo\ 1} and
\texttt{thermo\_modify\ format\ float\ \%20.15g}. The total energy
(etotal) was recorded at each step for both the serial and OMP-4-thread
runs of \texttt{born/coul/dsf} (864 atoms FCC) and \texttt{lj/cut/soft}
(864 atoms FCC with λ=0.5). In both cases:

\begin{enumerate}
\def\labelenumi{\arabic{enumi}.}
\tightlist
\item
  The total energy trajectories overlay exactly: the maximum point-wise
  difference \textbar E\_serial − E\_OMP4t\textbar{} across all 251
  steps is \textbf{8.97×10⁻¹⁴} (\texttt{born/coul/dsf}) and
  \textbf{4.11×10⁻¹³} (\texttt{lj/cut/soft}), corresponding to agreement
  at \textgreater13 decimal places.
\item
  Total energy drift over 250 steps: \textbf{0.021\%} for
  \texttt{born/coul/dsf} (charge-neutral ionic FCC lattice,
  well-equilibrated initial state). For \texttt{lj/cut/soft} at λ = 0.5
  the drift is \textbf{0.65\%}, dominated by a transient in the first
  \textasciitilde50 steps as the FCC crystal lattice relaxes into the
  soft-core potential well --- the trajectory stabilises around a mean
  after step 50, and the NVE integrator conserves the running total
  energy throughout. Both results confirm NVE stability of the generated
  OMP code.
\end{enumerate}

The bit-identical agreement confirms that the thread-safe force
accumulation (\texttt{thr-\textgreater{}get\_f()} +
\texttt{ev\_tally\_thr()} + \texttt{reduce\_thr()}) reproduces
floating-point sums in the same order as the serial implementation, as
required for scientific reproducibility.

\begin{center}\rule{0.5\linewidth}{0.5pt}\end{center}

\subsection{Notes on Key Technical
Decisions}\label{notes-on-key-technical-decisions}

\subsubsection{Why separate equilibration per lambda
window?}\label{why-separate-equilibration-per-lambda-window}

The standard practice in thermodynamic integration is to equilibrate
each λ window independently, rather than using sequential tempering
(starting from λ=0 and stepping to λ=1 sequentially). Independent
windows are preferred because: 1. They eliminate serial correlation
between windows, enabling independent error estimation 2. They guarantee
that each window samples its own equilibrium ensemble 3. They allow
trivial parallelization

The 40 ps production period per window is shorter than ideal for
production calculations (10--100 ns is recommended for accurate ΔG);
however, it is sufficient to demonstrate the methodology, to verify the
correct sign and ordering of ΔG values, and to show that the OMP
pipeline produces physically meaningful results. A full quantitative
comparison with experimental desolvation free energies would require
longer simulations and finite-size corrections (see S2.5).

\subsubsection{\texorpdfstring{Why \texttt{lj/cut/coul/long/soft} and
not
\texttt{lj/cut/soft}?}{Why lj/cut/coul/long/soft and not lj/cut/soft?}}\label{why-ljcutcoullongsoft-and-not-ljcutsoft}

EC, PC, and DME are polar solvents with permanent partial charges. The
dominant solvation contribution is Coulombic (Li⁺ charge +1 e attracting
O partial charges −0.35 to −0.50 e). Using
\texttt{lj/cut/coul/long/soft} simultaneously scales both LJ and Coulomb
interactions, which is physically correct for full annihilation of an
ion. Using only \texttt{lj/cut/soft} would leave the Coulomb interaction
at full strength, giving physically incorrect dU/dλ values and ΔG.

\subsubsection{Why Lorentz--Berthelot and not geometric
mixing?}\label{why-lorentzberthelot-and-not-geometric-mixing}

OPLS-AA specifies Lorentz--Berthelot (arithmetic σ, geometric ε) mixing
rules. These are enforced via \texttt{pair\_modify\ mix\ arithmetic} in
LAMMPS (arithmetic σ) combined with the geometric default for ε. Li⁺ was
parameterized in SPC/E water by Joung \& Cheatham (2008) using
Lorentz--Berthelot mixing; applying the same rules to OPLS-AA solvent
atoms is the standard practice in electrolyte FEP literature.

\begin{center}\rule{0.5\linewidth}{0.5pt}\end{center}

\subsection{Supplementary Table S6: Pair files modified by the
force-kernel
rewrite}\label{supplementary-table-s6-pair-files-modified-by-the-force-kernel-rewrite}

\emph{Table S6: the 225 pair files whose serial force kernels were
rewritten to use \texttt{\_\_restrict\_\_}-qualified typed pointers and
local accumulators, grouped by LAMMPS package. The complete file-level
inventory --- source path, package, \texttt{pair\_style} name recovered
from each file's registration macro, and which of the three code markers
(\texttt{\_noalias}, \texttt{fxtmp}, \texttt{dbl3\_t}) each file carries
--- is distributed with the software as
\texttt{Paper/tableS6\_a3\_files.csv}.}

{\def\LTcaptype{none} % do not increment counter
\begin{longtable}[]{@{}
  >{\raggedright\arraybackslash}p{(\linewidth - 4\tabcolsep) * \real{0.3333}}
  >{\raggedright\arraybackslash}p{(\linewidth - 4\tabcolsep) * \real{0.3333}}
  >{\raggedright\arraybackslash}p{(\linewidth - 4\tabcolsep) * \real{0.3333}}@{}}
\toprule\noalign{}
\begin{minipage}[b]{\linewidth}\raggedright
Package
\end{minipage} & \begin{minipage}[b]{\linewidth}\raggedright
Files
\end{minipage} & \begin{minipage}[b]{\linewidth}\raggedright
Representative styles
\end{minipage} \\
\midrule\noalign{}
\endhead
\bottomrule\noalign{}
\endlastfoot
EXTRA-PAIR & 42 & \texttt{beck}, \texttt{born/coul/dsf},
\texttt{born/coul/wolf} \\
MANYBODY & 20 & \texttt{adp}, \texttt{atm}, \texttt{comb3} \\
KSPACE & 19 & \texttt{born/coul/long}, \texttt{born/coul/msm},
\texttt{buck/coul/long} \\
(top level) & 14 & \texttt{born}, \texttt{buck}, \texttt{coul/cut} \\
FEP & 13 & \texttt{coul/cut/soft}, \texttt{coul/cut/soft/gapsys},
\texttt{coul/long/soft} \\
CG-DNA & 10 & \texttt{oxdna2/coaxstk}, \texttt{oxdna2/dh},
\texttt{oxdna3/xstk} \\
CORESHELL & 8 & \texttt{born/coul/dsf/cs}, \texttt{born/coul/long/cs},
\texttt{born/coul/wolf/cs} \\
DIELECTRIC & 7 & \texttt{coul/cut/dielectric},
\texttt{coul/long/dielectric}, \texttt{lj/cut/coul/cut/dielectric} \\
INTERLAYER & 7 & \texttt{coul/shield}, \texttt{ilp/graphene/hbn},
\texttt{ilp/tmd} \\
MOLECULE & 7 & \texttt{hbond/dreiding/lj},
\texttt{hbond/dreiding/morse}, \texttt{lj/charmm/coul/charmm} \\
SPIN & 7 & \texttt{spin/dipole/cut}, \texttt{spin/dipole/long},
\texttt{spin/dmi} \\
DPD-REACT & 6 & \texttt{dpd/fdt}, \texttt{dpd/fdt/energy},
\texttt{exp6/rx} \\
SPH & 6 & \texttt{sph/heatconduction}, \texttt{sph/idealgas},
\texttt{sph/lj} \\
ASPHERE & 5 & \texttt{gayberne}, \texttt{line/lj}, \texttt{resquared} \\
DPD-BASIC & 5 & \texttt{dpd}, \texttt{dpd/coul/slater/long},
\texttt{dpd/ext} \\
PERI & 5 & \texttt{peri/eps}, \texttt{peri/lps}, \texttt{peri/pmb} \\
COLLOID & 4 & \texttt{brownian}, \texttt{brownian/poly},
\texttt{colloid} \\
DPD-MESO & 4 & \texttt{edpd}, \texttt{mdpd}, \texttt{mdpd/rhosum} \\
CLASS2 & 3 & \texttt{lj/class2}, \texttt{lj/class2/coul/cut},
\texttt{lj/class2/coul/long} \\
DIPOLE & 3 & \texttt{lj/cut/dipole/cut}, \texttt{lj/cut/dipole/long},
\texttt{lj/sf/dipole/sf} \\
DRUDE & 3 & \texttt{coul/tt}, \texttt{lj/cut/thole/long},
\texttt{thole} \\
GRANULAR & 3 & \texttt{gran/hertz/history}, \texttt{gran/hooke},
\texttt{gran/hooke/history} \\
SMTBQ & 3 & \texttt{smatb}, \texttt{smatb/single}, \texttt{smtbq} \\
APIP & 2 & \texttt{eam/apip}, \texttt{lambda/input/csp/apip} \\
MACHDYN & 2 & \texttt{smd/hertz}, \texttt{smd/tri\_surface} \\
MISC & 2 & \texttt{agni}, \texttt{tracker} \\
MOFFF & 2 & \texttt{buck6d/coul/gauss/dsf},
\texttt{buck6d/coul/gauss/long} \\
RHEO & 2 & \texttt{rheo}, \texttt{rheo/solid} \\
YAFF & 2 & \texttt{lj/switch3/coulgauss/long},
\texttt{mm3/switch3/coulgauss/long} \\
AMOEBA & 1 & \texttt{amoeba} \\
BODY & 1 & \texttt{body/nparticle} \\
BPM & 1 & \texttt{bpm/spring} \\
DPD-SMOOTH & 1 & \texttt{sdpd/taitwater/isothermal} \\
EFF & 1 & \texttt{eff/cut} \\
MC & 1 & \texttt{dsmc} \\
ML-SNAP & 1 & \texttt{snap} \\
ML-UF3 & 1 & \texttt{uf3} \\
PYTHON & 1 & \texttt{python} \\
\textbf{Total} & \textbf{225} & \\
\end{longtable}
}

Three further files were rewritten and then reverted, because the
rewrite did not preserve floating-point operation order for them:
\texttt{pair\_buck\_coul\_cut}, \texttt{pair\_kolmogorov\_crespi\_z},
and \texttt{pair\_polymorphic} (see Correctness Verification). They are
counted in neither the table above nor the 225 total.

Styles left unmodified are those without a standard half-list or
full-list neighbor loop. Contrary to our initial expectation, multi-body
potentials (EAM, Tersoff, SW, Vashishta) and two machine-learning
potentials (SNAP, UF3) proved amenable to the transformation and are
included above; all of them reproduce the reference forces and energies
of the unmodified code.

\begin{center}\rule{0.5\linewidth}{0.5pt}\end{center}

\subsection{Supplementary Table S7: A-3 Before/After Performance
Benchmarks}\label{supplementary-table-s7-a-3-beforeafter-performance-benchmarks}

\emph{Table S7: Loop-time and pair-time improvement for three
representative pair styles after A-3 backporting. Benchmarks: 864-atom
FCC cell, 500 NVE steps, serial 1T, N=3 runs, mean ± std. Loop
improvement = (before − after)/before × 100\%; pair improvement uses
pair time column.}

{\def\LTcaptype{none} % do not increment counter
\begin{longtable}[]{@{}
  >{\raggedright\arraybackslash}p{(\linewidth - 12\tabcolsep) * \real{0.1429}}
  >{\raggedright\arraybackslash}p{(\linewidth - 12\tabcolsep) * \real{0.1429}}
  >{\raggedright\arraybackslash}p{(\linewidth - 12\tabcolsep) * \real{0.1429}}
  >{\raggedright\arraybackslash}p{(\linewidth - 12\tabcolsep) * \real{0.1429}}
  >{\raggedright\arraybackslash}p{(\linewidth - 12\tabcolsep) * \real{0.1429}}
  >{\raggedright\arraybackslash}p{(\linewidth - 12\tabcolsep) * \real{0.1429}}
  >{\raggedright\arraybackslash}p{(\linewidth - 12\tabcolsep) * \real{0.1429}}@{}}
\toprule\noalign{}
\begin{minipage}[b]{\linewidth}\raggedright
Style
\end{minipage} & \begin{minipage}[b]{\linewidth}\raggedright
Before loop (s)
\end{minipage} & \begin{minipage}[b]{\linewidth}\raggedright
After loop (s)
\end{minipage} & \begin{minipage}[b]{\linewidth}\raggedright
Loop improvement
\end{minipage} & \begin{minipage}[b]{\linewidth}\raggedright
Before pair (s)
\end{minipage} & \begin{minipage}[b]{\linewidth}\raggedright
After pair (s)
\end{minipage} & \begin{minipage}[b]{\linewidth}\raggedright
Pair improvement
\end{minipage} \\
\midrule\noalign{}
\endhead
\bottomrule\noalign{}
\endlastfoot
momb & 1.768 ± 0.018 & 1.643 ± 0.032 & \textbf{+7.1\%} & 1.706 & 1.579 &
+7.5\% \\
coul/ctip & 22.25 ± 1.06 & 20.97 ± 0.54 & \textbf{+5.7\%} & 22.17 &
20.90 & +5.7\% \\
dispersion/d3 & 15.31 ± 0.35 & 15.81 ± 1.95 & ≈0\% (noise) & 15.27 &
15.76 & --- \\
\end{longtable}
}

The \texttt{dispersion/d3} style uses a multi-pass loop structure that
already has loop-carried dependencies from three-body dispersion; the
restrict/fxtmp pattern provides no benefit in this case and was left
unchanged in the final submission.

\begin{center}\rule{0.5\linewidth}{0.5pt}\end{center}

\subsection{Supplementary Figure S3: A-3 Code
Transformation}\label{supplementary-figure-s3-a-3-code-transformation}

\emph{Figure S3: Representative before/after code transformation applied
in Step A-3 (compiler-friendly serial hotloop backporting). The pattern
was applied to \textasciitilde102 of 307 pair\_style files.}

\textbf{Before (standard \texttt{pair\_lj\_cut.cpp} hotloop):}

\begin{Shaded}
\begin{Highlighting}[]
\DataTypeTok{double} \OperatorTok{**}\NormalTok{x }\OperatorTok{=}\NormalTok{ atom}\OperatorTok{{-}\textgreater{}}\NormalTok{x}\OperatorTok{;}
\DataTypeTok{double} \OperatorTok{**}\NormalTok{f }\OperatorTok{=}\NormalTok{ atom}\OperatorTok{{-}\textgreater{}}\NormalTok{f}\OperatorTok{;}

\ControlFlowTok{for} \OperatorTok{(}\NormalTok{ii }\OperatorTok{=} \DecValTok{0}\OperatorTok{;}\NormalTok{ ii }\OperatorTok{\textless{}}\NormalTok{ inum}\OperatorTok{;}\NormalTok{ ii}\OperatorTok{++)} \OperatorTok{\{}
\NormalTok{  i }\OperatorTok{=}\NormalTok{ ilist}\OperatorTok{[}\NormalTok{ii}\OperatorTok{];}
  \CommentTok{// ... inner loop ...}
  \ControlFlowTok{for} \OperatorTok{(}\NormalTok{jj }\OperatorTok{=} \DecValTok{0}\OperatorTok{;}\NormalTok{ jj }\OperatorTok{\textless{}}\NormalTok{ jnum}\OperatorTok{;}\NormalTok{ jj}\OperatorTok{++)} \OperatorTok{\{}
    \CommentTok{// ...}
\NormalTok{    f}\OperatorTok{[}\NormalTok{i}\OperatorTok{][}\DecValTok{0}\OperatorTok{]} \OperatorTok{+=}\NormalTok{ delx }\OperatorTok{*}\NormalTok{ fpair}\OperatorTok{;}   \CommentTok{// per{-}pair write to f[i]: loop{-}carried dependency}
\NormalTok{    f}\OperatorTok{[}\NormalTok{i}\OperatorTok{][}\DecValTok{1}\OperatorTok{]} \OperatorTok{+=}\NormalTok{ dely }\OperatorTok{*}\NormalTok{ fpair}\OperatorTok{;}
\NormalTok{    f}\OperatorTok{[}\NormalTok{i}\OperatorTok{][}\DecValTok{2}\OperatorTok{]} \OperatorTok{+=}\NormalTok{ delz }\OperatorTok{*}\NormalTok{ fpair}\OperatorTok{;}
    \ControlFlowTok{if} \OperatorTok{(}\NormalTok{newton\_pair }\OperatorTok{||}\NormalTok{ j }\OperatorTok{\textless{}}\NormalTok{ nlocal}\OperatorTok{)} \OperatorTok{\{}
\NormalTok{      f}\OperatorTok{[}\NormalTok{j}\OperatorTok{][}\DecValTok{0}\OperatorTok{]} \OperatorTok{{-}=}\NormalTok{ delx }\OperatorTok{*}\NormalTok{ fpair}\OperatorTok{;}
\NormalTok{      f}\OperatorTok{[}\NormalTok{j}\OperatorTok{][}\DecValTok{1}\OperatorTok{]} \OperatorTok{{-}=}\NormalTok{ dely }\OperatorTok{*}\NormalTok{ fpair}\OperatorTok{;}
\NormalTok{      f}\OperatorTok{[}\NormalTok{j}\OperatorTok{][}\DecValTok{2}\OperatorTok{]} \OperatorTok{{-}=}\NormalTok{ delz }\OperatorTok{*}\NormalTok{ fpair}\OperatorTok{;}
    \OperatorTok{\}}
  \OperatorTok{\}}
\OperatorTok{\}}
\end{Highlighting}
\end{Shaded}

\textbf{After (A-3 backported pattern):}

\begin{Shaded}
\begin{Highlighting}[]
\AttributeTok{const} \KeywordTok{auto} \OperatorTok{*}\NormalTok{ \_noalias }\AttributeTok{const}\NormalTok{ x }\OperatorTok{=} \OperatorTok{(}\DataTypeTok{dbl3\_t} \OperatorTok{*)}\NormalTok{ atom}\OperatorTok{{-}\textgreater{}}\NormalTok{x}\OperatorTok{[}\DecValTok{0}\OperatorTok{];}  \CommentTok{// restrict: no aliasing}
\KeywordTok{auto} \OperatorTok{*}\NormalTok{ \_noalias }\AttributeTok{const}\NormalTok{ f }\OperatorTok{=} \OperatorTok{(}\DataTypeTok{dbl3\_t} \OperatorTok{*)}\NormalTok{ atom}\OperatorTok{{-}\textgreater{}}\NormalTok{f}\OperatorTok{[}\DecValTok{0}\OperatorTok{];}

\ControlFlowTok{for} \OperatorTok{(}\NormalTok{ii }\OperatorTok{=} \DecValTok{0}\OperatorTok{;}\NormalTok{ ii }\OperatorTok{\textless{}}\NormalTok{ inum}\OperatorTok{;}\NormalTok{ ii}\OperatorTok{++)} \OperatorTok{\{}
\NormalTok{  i }\OperatorTok{=}\NormalTok{ ilist}\OperatorTok{[}\NormalTok{ii}\OperatorTok{];}
  \DataTypeTok{double}\NormalTok{ fxtmp }\OperatorTok{=} \FloatTok{0.0}\OperatorTok{,}\NormalTok{ fytmp }\OperatorTok{=} \FloatTok{0.0}\OperatorTok{,}\NormalTok{ fztmp }\OperatorTok{=} \FloatTok{0.0}\OperatorTok{;}         \CommentTok{// local accumulators}

  \ControlFlowTok{for} \OperatorTok{(}\NormalTok{jj }\OperatorTok{=} \DecValTok{0}\OperatorTok{;}\NormalTok{ jj }\OperatorTok{\textless{}}\NormalTok{ jnum}\OperatorTok{;}\NormalTok{ jj}\OperatorTok{++)} \OperatorTok{\{}
    \CommentTok{// ...}
\NormalTok{    fxtmp }\OperatorTok{+=}\NormalTok{ delx }\OperatorTok{*}\NormalTok{ fpair}\OperatorTok{;}   \CommentTok{// accumulate locally: no memory dependency}
\NormalTok{    fytmp }\OperatorTok{+=}\NormalTok{ dely }\OperatorTok{*}\NormalTok{ fpair}\OperatorTok{;}
\NormalTok{    fztmp }\OperatorTok{+=}\NormalTok{ delz }\OperatorTok{*}\NormalTok{ fpair}\OperatorTok{;}
    \ControlFlowTok{if} \OperatorTok{(}\NormalTok{newton\_pair }\OperatorTok{||}\NormalTok{ j }\OperatorTok{\textless{}}\NormalTok{ nlocal}\OperatorTok{)} \OperatorTok{\{}
\NormalTok{      f}\OperatorTok{[}\NormalTok{j}\OperatorTok{].}\NormalTok{x }\OperatorTok{{-}=}\NormalTok{ delx }\OperatorTok{*}\NormalTok{ fpair}\OperatorTok{;}
\NormalTok{      f}\OperatorTok{[}\NormalTok{j}\OperatorTok{].}\NormalTok{y }\OperatorTok{{-}=}\NormalTok{ dely }\OperatorTok{*}\NormalTok{ fpair}\OperatorTok{;}
\NormalTok{      f}\OperatorTok{[}\NormalTok{j}\OperatorTok{].}\NormalTok{z }\OperatorTok{{-}=}\NormalTok{ delz }\OperatorTok{*}\NormalTok{ fpair}\OperatorTok{;}
    \OperatorTok{\}}
  \OperatorTok{\}}
\NormalTok{  f}\OperatorTok{[}\NormalTok{i}\OperatorTok{].}\NormalTok{x }\OperatorTok{+=}\NormalTok{ fxtmp}\OperatorTok{;}   \CommentTok{// single write per outer atom}
\NormalTok{  f}\OperatorTok{[}\NormalTok{i}\OperatorTok{].}\NormalTok{y }\OperatorTok{+=}\NormalTok{ fytmp}\OperatorTok{;}
\NormalTok{  f}\OperatorTok{[}\NormalTok{i}\OperatorTok{].}\NormalTok{z }\OperatorTok{+=}\NormalTok{ fztmp}\OperatorTok{;}
\OperatorTok{\}}
\end{Highlighting}
\end{Shaded}

\textbf{Key changes}: 1. \texttt{double\ **x/f} →
\texttt{dbl3\_t\ *\ \_noalias} cast: informs compiler that x and f do
not alias each other or atom type arrays, enabling auto-vectorization of
distance calculations. 2. Per-pair \texttt{f{[}i{]}{[}k{]}\ +=\ ...} →
post-inner-loop \texttt{f{[}i{]}.k\ +=\ ftmp}: eliminates loop-carried
memory dependency on the outer atom's force accumulation, allowing the
compiler to keep the running sum in a register across inner-loop
iterations. 3. \texttt{cutsq{[}itype{]}} and other 2D arrays cached as
1D \texttt{const\ double\ *} row pointers before the inner loop.

Floating-point operation order is preserved exactly (verified
bit-identical for all 102 modified files using
\texttt{compare\_thermo.py}).

\begin{center}\rule{0.5\linewidth}{0.5pt}\end{center}

\emph{End of Supplementary Materials}

\end{document}
